\begin{abstract}
\noindent
Các hệ thống sinh tăng cường truy hồi (Retrieval-Augmented Generation, RAG)
có khả năng cung cấp thông tin từ nguồn tri thức bên ngoài, nhưng vẫn gặp khó
khăn đối với các câu hỏi pháp luật yêu cầu kết nối nhiều quy tắc, xác định
đường suy luận và kiểm tra các giả định phản thực tế. Nghiên cứu này trình
bày một pipeline CausalRAG cho hỏi--đáp pháp luật hình sự Việt Nam, trong đó
các quy tắc pháp luật được chuẩn hóa thành quan hệ điều kiện--hệ quả và tổ
chức dưới dạng đồ thị nhân quả dị thể. Hệ thống kết hợp bộ nhớ vector sử dụng
BGE-M3, truy hồi đường dẫn nhân quả nhiều bước, xác minh claim theo nội dung
truy vấn và sinh câu trả lời có căn cứ điều luật.

Thành phần xác minh trung tâm của hệ thống là mô hình nhân quả cấu trúc pháp
lý (Legal Structural Causal Model, LegalSCM), một mô hình suy luận quy phạm
tất định với miền trạng thái ba giá trị. LegalSCM xác nhận factual causal
path và thực hiện hard intervention bằng cách gán sự kiện trung gian về
trạng thái sai, cắt các cơ chế đầu vào tương ứng và suy luận lại outcome.
Một baseline dựa trên xóa node và kiểm tra khả năng tiếp cận trên đồ thị
được sử dụng để thực hiện paired ablation.

Thực nghiệm được tiến hành trên 2.884 quy tắc, 1.814 sự kiện và một bộ đánh
giá tổng hợp gồm 250 câu hỏi đa bước. Hệ thống đạt Rule Recall@5 là 69,93\%,
Event Recall@5 là 74,66\%, độ chính xác xác minh 92,00\%, Answer Token F1 là
68,69\% và Citation F1 là 62,86\%. Top-1 Exact Path đạt 35,20\%, trong khi
Oracle Exact Path đạt 71,20\%, cho thấy xếp hạng đường dẫn vẫn là nút thắt
chính. Paired ablation không ghi nhận khác biệt về chất lượng giữa
LegalSCM và path ablation trên benchmark hiện tại; tuy nhiên, LegalSCM cung
cấp semantics can thiệp và trace suy luận chi tiết hơn với thời gian xử lý
trung bình cao hơn khoảng 3,23 lần. Kết quả cho thấy tiềm năng của việc kết
hợp truy hồi theo đồ thị, xác minh có cấu trúc và sinh câu trả lời có căn cứ,
đồng thời chỉ ra nhu cầu xây dựng benchmark phản thực tế nhạy hơn với cơ chế
can thiệp.
\end{abstract}

\vspace{0.35em}
\noindent\textbf{Từ khóa:}
Retrieval-Augmented Generation; CausalRAG; đồ thị nhân quả;
mô hình nhân quả cấu trúc; suy luận phản thực tế;
hỏi--đáp pháp luật; Bộ luật Hình sự Việt Nam.
\vspace{0.6em}
